\section{Limitations}
\label{app:limitations}
Despite UBTree's strong empirical performance, two limitations remain.

\noindent\textit{Tree drafting is sensitive to very high concurrency.}
Speculative decoding is motivated by the spare compute available when autoregressive inference is limited by memory bandwidth, as is common in LLM decoding. Verifying multiple draft tokens together uses this capacity to amortize memory-access costs. As concurrency increases and inference becomes compute-bound, the spare capacity available for speculation shrinks, potentially eliminating speedups for both single-path and tree-based methods. Sensitivity to this transition is shared by all tree-based drafters and is not specific to UBTree. Tree verification spends additional compute on alternative continuations and therefore favors the part of the compute--bandwidth spectrum with more spare compute. Adaptive tree drafting can mitigate this limitation by varying the verification budget with available compute: branching allows the number of verified candidates to vary beyond the token count of a fixed single-path block, although the proposer block size still bounds draft depth. Our load-dependent tree budgets already retain speedups at the tested concurrency levels (Table~\ref{tab:concurrency-results}), and more flexible runtime adaptation could further improve resource use. UBTree can also serve low-concurrency workloads where spare compute is abundant. We therefore view this limitation as a constraint on the favorable deployment regime, rather than a fundamental obstacle to UBTree's practical value.

\noindent\textit{Adaptive tree strategies rely on searched configurations.}
Our current adaptive strategy uses tree budgets $B$ and search widths $W$ selected through configuration search for each model and concurrency level. These settings are fixed for a given serving load; they do not determine the budget from the transition-score distribution of each decoding round. Although the tree's contents depend on the predicted scores, the budget-selection policy remains empirically tuned and is not guaranteed to be optimal as prediction uncertainty and available compute change. A more principled strategy would select tree sizes at each decoding round using calibrated estimates of acceptance benefit derived from transition scores, together with a cost model reflecting hardware and serving load. We leave developing and evaluating such a policy to future work.

\section{Additional Related Work}
\label{app:additional-related-work}

\noindent\textbf{Speculative Decoding.} In addition to autoregressive~\citep{leviathan2023fast,chen2023accelerating,li2024eagle}, parallel~\citep{chen2026dflash,liu2026dart,inco2026dflash2}, semi-autoregressive~\citep{huang2026domino,cheng2026dspark}, and tree-based~\citep{cai2024medusa,miao2024specinfer,li2024eagle2,li2026eagle,ringel2026accelerating,li2026dartree} drafting methods mentioned, existing research also explored mixing methods for speculative decoding. PARD adapts autoregressive models for parallel token prediction~\citep{an2025pard}, while DiffuSpec uses pretrained diffusion language models as drafters~\citep{li2026diffuspec}.

\noindent\textbf{Lossy Acceleration for Large Language Models.}
Lossy acceleration reduces inference cost by allowing deviations from the target model's output distribution. BiLD~\citep{kim2023biglittle} coordinates small and large models through confidence-based fallback and discrepancy-based rollback. The lossy extension of DistillSpec~\citep{zhou2024distillspec} relaxes token acceptance using a lenience factor, while Fuzzy Speculative Decoding~\citep{holsman2025fuzzy} bases acceptance on divergence between draft and target distributions. Speculative cascades~\citep{narasimhan2025faster} implement model-deferral rules through speculative execution, selecting between draft and target distributions; DIVERSED~\citep{wang2026diversed} instead learns context-dependent mixtures of these distributions for verification. Judge Decoding~\citep{bachmann2025judge,garipov2026autojudge} accepts otherwise-rejected tokens using a lightweight correctness classifier on target hidden states. In contrast to these relaxed verification methods, UBTree improves candidate coverage and path quality through joint proposer--selector training while retaining target-distribution-preserving verification.

\noindent\textbf{Multi-token Language Models for Fast Generation.}
Another line of work incorporates multi-token prediction into the language model itself. \citet{gloeckle2024better} train multiple prediction heads on a shared backbone to predict several future tokens simultaneously. Diffusion language models, including LLaDA~\citep{nie2025large} and Dream~\citep{xie2025dream}, learn to recover masked tokens and generate text through iterative parallel denoising. LLaDA is trained from scratch, whereas Dream adapts pretrained autoregressive weights. Block diffusion~\citep{arriola2025block} combines autoregressive generation across blocks with parallel denoising within each block, enabling flexible-length generation and prefix KV caching. Large-scale block diffusion with different recipes enriches the explored space of taming diffusion for multi-token fast generation, including WeDLM~\citep{liu2025wedlm}, Mercury~\citep{inception2025mercury}, and DiffusionGemma~\citep{diffusiongemma2026report}. These approaches modify the generative model or its training to enable multi-token generation, while having the risk of degrading the sample quality.
